\section{Codec Design}
\label{sec:codecs}

Both codecs share the same high-level pipeline: read the raw big-endian
16-bit image; apply a spatial predictor; zigzag-encode the signed residual
$r = \text{pixel} - \text{pred}$ as an unsigned symbol $z = 2r$ if $r\ge0$,
$z = -2r-1$ otherwise; entropy-code the result.

% ── prism-block ───────────────────────────────────────────────────────────────
\subsection{prism-block: Speed-Focused}
\label{sec:speed}

\paragraph{Image partitioning.}
The image is split into horizontal bands of 150 rows. Bands are compressed
independently in parallel, with each thread atomically claiming the next
unclaimed band.

\paragraph{Predictor.}
Each pixel uses the \emph{Gradient-Adjusted Prediction} (GAP) function of
its three causal neighbours: left ($W$), above ($N$), and above-left ($NW$).
GAP estimates directional gradients $d_h = |W - NW|$ and $d_v = |N - NW|$.
If the horizontal gradient dominates ($d_h > 2\,d_v$), the pixel above is
used; if the vertical gradient dominates, the pixel to the left is used;
otherwise a weighted blend is computed:
\[
  \hat{p} = \operatorname{clamp}\!\left(\frac{(d_v+1)\,W + (d_h+1)\,N}{d_v+d_h+2},\;\min(W,N),\;\max(W,N)\right).
\]

\paragraph{Entropy coding.}
After zigzag encoding, the 16-bit symbol $z$ is split into a high byte
$z_h = z \gg 8$ and a low byte $z_l = z\,\&\,0\!\mathrm{xFF}$.
Both are coded with a \emph{range coder} (carry-based, 64-bit state) backed
by adaptive Fenwick-tree frequency models.
$z_h$ uses a single shared model; $z_l$ uses one of eight models selected by
the context
\[
  c = \operatorname{bin}(|\hat{r}_W| + |\hat{r}_N|),
\]
where $\hat{r}_W, \hat{r}_N$ are the residuals of the left and upper pixel and
the bins span the ranges $\{0\},\{1\text{--}2\},\{3\text{--}8\},\{9\text{--}32\},
\{33\text{--}64\},\{65\text{--}128\},\{129\text{--}512\},\{>512\}$.
Each model halves all counts when the total reaches $2^{16}$, tracking
distribution shifts within a band.

% ── HAIS ──────────────────────────────────────────────────────────────────────
\subsection{HAIS: Balanced}
\label{sec:balanced}

\subsubsection*{Block Selection}
HAIS infers image dimensions from the file size (attempting the integer square
root, then a list of common astronomical widths: 1500, 2048, 4096, \ldots).
It then selects the largest block side $b\le512$ that divides both $W$ and $H$
exactly while producing at least four blocks. For $1500\times1500$, this gives
$b=500$ and a $3\times3$ grid of nine blocks.

\subsubsection*{Predictive Modes}
For each block, HAIS evaluates five candidate predictors and picks the one with
the lowest estimated entropy cost.

\begin{description}[leftmargin=2em, labelindent=0em]
  \item[Mode~0 (raw)] Symbols are the raw pixel values; no prediction applied.
  \item[Mode~1 (avg)] Prediction is the three-neighbour average
    $\hat{p} = \lfloor(W + N + NW + 1)/3\rfloor$.
  \item[Mode~3 (gmean)] Prediction is the global mean $\bar{g}$ of the entire
    image, computed once before blocking. This single subtraction removes the
    bias pedestal common in all frames.
  \item[Mode~4 (LS-4)] Ordinary least-squares with four features
    $(W, N, NW, 1)$, fit per block via Gauss-Jordan elimination.
    The four float32 weights (16\,B) are stored in the block header.
  \item[Mode~6 (LS-5)] Extends Mode~4 with north-east ($NE$) and two-rows-above
    ($NN$) neighbours: features $(W, N, NW, NE, NN)$.
    Five float32 weights (20\,B) are stored in the block header.
\end{description}

Both LS modes derive their weights from the block itself — no training data or
separate pass is required. The weight overhead is included in the mode
selection cost:
\[
  \text{cost}_m = H(\text{hi}_8) + H(\text{lo}_8) + \frac{8\,\delta_m}{n_\text{pix}},
\]
where $H(\cdot)$ is the empirical byte entropy, $n_\text{pix}$ is the block
pixel count, and $\delta_m \in \{0,0,0,16,20\}$ is the weight header size.

\subsubsection*{Entropy Coding: FSE/tANS}
Residuals are split into a high-byte stream (\texttt{hi}) and a low-byte
stream (\texttt{lo}), each coded independently with Finite State
Entropy~\cite{collet2013fse}.

The FSE table has $\text{SCALE}=2^{16}=65536$ states. Symbol frequencies are
normalised to sum to SCALE and spread across states using Collet's
deterministic formula with step $= (\text{SCALE}\gg1)+(\text{SCALE}\gg3)+3$.
The encoder processes symbols in reverse order, emitting variable-length
transition bits; the decoder reads the two-byte initial state and reconstructs
symbols in the forward order.

The compressed stream header encodes the frequency table in one of three
formats: \emph{single-symbol} (\texttt{0x01}, 2\,B total) when all values are
identical; \emph{sparse} (\texttt{nnz} as flag byte followed by
$\text{nnz}\times5$\,B symbol--freq pairs) for up to 254 distinct symbols;
or \emph{dense} (\texttt{0x00} followed by $256\times4$\,B) otherwise.

\subsubsection*{Context FSE}
When the fraction of zero high bytes in a block falls between 2\% and 98\%,
HAIS additionally tries splitting \texttt{lo} into two conditional streams:
$\texttt{lo\_zero}$ (lo bytes where $\text{hi}=0$) and $\texttt{lo\_nonzero}$
(all others). The three-stream payload is used only if it is strictly smaller
than the two-stream baseline. The high-byte stream is encoded once and reused
in both variants, avoiding redundant computation. Bit~7 of the mode byte
signals to the decompressor that the three-stream layout is active.

\subsubsection*{Parallel Execution}
Compression and decompression both use a thread pool backed by a single
\texttt{std::atomic<int>} counter. Each thread atomically increments the
counter with \texttt{fetch\_add} to claim the next block; no mutex is required
because block results are written to pre-allocated slots.
